% !TeX encoding = UTF-8
\chapter{故障排查与安全加固}

软件系统上线运行后，会面临复杂的宿主机物理环境变化、外部网络波动以及恶意攻击威胁。为了最大程度降低故障平均修复时间并增强系统在公网环境下的抗风险能力，本章将依据计算机软件文档编制规范~\cite{gb8567} 的系统安全与健壮性设计要求，梳理 LingoBridge 平台部署期间的常见故障诊断矩阵，提供针对生产主机的安全加固建议，并为后续系统从最小可行性产品平滑演进至大规模高可用架构提出具体的改造设计方案。

\section{常见故障排查}

在平台部署或运行期间，若系统出现前端资源无法加载、接口请求报错或业务逻辑中断，运维人员需要根据故障表现症状，利用系统命令对关键瓶颈进行快速定位。下表归纳了系统部署验证中常见的四大典型故障及其对应的潜在诱发原因和推荐的排除解决方法。

\begin{table}[htbp]
  \centering
  \caption{LingoBridge 部署常见故障诊断与排除方案对照表}
  \label{tab:troubleshooting-matrix}
  \begin{tabular}{lll}
    \toprule
    故障表现症状 & 潜在诱发原因 & 推荐诊断与排除解决方法 \\
    \midrule
    服务启动失败且提示端口冲突 & 服务器中其他进程占用了 \num{3001} 或 \num{80} 端口 & 执行端口占用检查命令并杀掉冲突进程以释放资源 \\
    用户登录或注册提示数据库写入错误 & 本地数据库文件或数据目录无写权限 & 运行权限授予命令赋予当前执行用户读写权限 \\
    浏览器请求时提示 502 错误 & 后端接口服务崩溃或内部代理转发端口配置错误 & 检查后端进程运行状态并重新对齐代理配置文件 \allowbreak \\
    前端静态资源页面加载报 404 错误 & 前端静态打包产物目录缺失或代理路径指定有误 & 验证前端产物目录完整性并更新代理配置根目录 \\
    \bottomrule
  \end{tabular}
\end{table}

第一类是服务端口占用故障。如果系统启动失败并抛出监听端口已被绑定的网络异常，运维人员需要使用终端命令 \texttt{lsof -i :3001} 或 \texttt{netstat -tuln} 查询当前占用三千零一端口的进程识别号，在确认其不是核心系统依赖后，可通过强制终止命令将其杀掉，以释放端口资源供后端 API 服务正常绑定。

第二类是本地数据库写入故障。当发生新用户无法注册或跟读数据无法持久化，且控制台报错提示磁盘写入受阻时，表明 \texttt{backend/data/db.json} 数据文件或其所在目录的读写属性发生了错乱。运维人员需在宿主机运行 \texttt{chmod -R 755 backend/data} 指令，并使用所有权变更指令 \texttt{chown} 将该文件夹的归属权重新划归给当前 Node 进程的执行用户，从而恢复文件数据库的正常读写。

第三类是反向代理网关五零二故障。如果浏览器访问页面时弹出该错误，意味着反向代理服务器已启动但无法建立与后端应用服务器的连接。此时，需要登录云主机，先查看后端 Express 进程是否处于活跃状态，再核查反向代理配置文件中的代理转发目标端口是否与后端监听端口完全一致。如果处于 Docker 容器网络中，还必须确保后端容器与代理容器加入了相同的虚拟子网，且代理转发目标配置为了正确的后端容器主机名。

第四类是前端静态资源加载四零四故障。如果页面显示空白且控制台充满资源缺失的四零四报错，一般是因为前端打包静态产物目录没有被正确挂载，或者是在前端项目构建时配置的基准路径与代理配置中的静态根目录发生了漂移。运维人员需仔细对齐 Vite 配置文件中的基准路径参数，并确认编译生成的 \texttt{dist} 目录已完整输出并被反向代理的静态根路径所指向。

\section{系统安全加固建议}

系统在公网正式开通后，直接暴露在无边界的互联网威胁之下，必须立即执行安全加固操作，以避免服务器沦为网络木马的攻击目标，保护师生的隐私多媒体数据。

首先，在物理与逻辑层级实施严格的端口访问控制。除了用于远程连接的安全外壳协议端口二十二，以及面向师生浏览器提供服务的超文本传输协议端口八十与四百四十三外，应当利用操作系统的网络防火墙或者云厂商提供的虚拟安全组，彻底屏蔽其他非必要物理端口的公网入站流量。特别需要指出的是，后端 API 监听的三千零一端口应当被严格限制为仅允许本地环回接口访问，杜绝任何外部网络直接探测或建立非代理连接的企图。

其次，强制推行全链路传输层安全协议加密。部署人员应从权威机构申请免费的数字证书，并将其配置文件挂载至反向代理中，在配置规则中明确指定仅启用高安全强度的加密套件，并将所有传统的明文超文本传输协议请求强制重定向为加密超文本传输协议请求，确保包括用户登录口令、跟读语音流在内的所有会话数据在跨国公共网络传输过程中不被非法窃听与篡改。

最后，遵循最小权限原则硬化宿主机文件系统。整个 LingoBridge 应用部署目录 \texttt{/opt/lingobridge} 的所有者，应当被设置为系统底层的非特权普通用户，严禁直接使用系统最高权限的超级用户来运行 Node 进程或拉起 Docker 容器，防止黑客通过应用层级的潜在漏洞提权控制整个宿主机系统。同时，对文件型数据库目录进行严格的访问控制，仅向应用进程账户开放读写权限。

\section{生产环境演进建议}

当前的 LingoBridge 平台完全基于轻量级的文件数据库和本地磁盘文件存放机制构建，此种架构能够支撑数十人规模的最小可行性测试，但无法满足未来更大规模用户并发、海量语音数据高频吞吐及服务高可用性的长远诉求。为此，系统应当按照如下方向进行高可用演进改造。

在状态存储层面，应当将本地文件数据库 \texttt{backend/data/db.json} 平滑迁移至成熟的关系型数据库管理系统，如开源的高性能对象关系数据库 PostgreSQL。在后端 Node 架构中，引入先进的对象关系映射器如 Prisma 或 TypeORM，这不仅能够大幅提升复杂业务查询的性能，还可以利用关系型数据库自带的行级锁和事务并发机制，彻底消除高并发写入下可能发生的文件锁死和数据覆盖故障。

在多媒体附件存储层面，应当将本地磁盘路径 \texttt{backend/storage} 升级为云服务商提供的分布式对象存储服务，如腾讯云的对象存储服务或亚马逊的简单存储服务。将教师端上传的大文件课件以及学生端提交的语音录音，通过对象存储服务专有的应用程序接口直接上传至分布式的对象存储桶中，不仅能够实现多节点并发的高速分发，还能利用云厂商的多副本灾备机制保障多媒体资产的零丢失。

在性能与高可用层面，应当在系统架构中引入 Redis 内存键值数据库，将其用于缓存高频访问的课件索引和用户会话数据，以减轻底层关系数据库的压力。此外，可以部署容器集群编排引擎将前后端服务进行水平多实例扩展，在前端设置负载均衡器进行请求的轮询分发，当某一服务节点发生物理故障时可由引擎自动进行健康检测并进行零感知无缝拉起，实现平台真正意义上的高并发、高可用与弹性伸缩。
